% 数制和码制 子文件，匹配PPT章节顺序，无手动编号，无冗余应试文字
\chapter{数制和码制}
\section{概述}
\begin{enumerate}
    \item 数制：带位权的进位计数规则；码制：信息向二进制映射的编码规则。
    \item 数字信号仅用0、1表征两种对立电平/开关状态，依靠数制、码制赋予数值与字符含义。
\end{enumerate}

\section{几种常用的数制}
\subsection{通用定义与位权展开公式}
基数$N$：数制内可用数码总个数；位权：$N^i$；任意$n$位整数、$m$位小数的$N$进制数展开：
\[
(X)_N=\sum_{i=-m}^{n-1}k_i \cdot N^i,\quad 0\le k_i\le N-1
\]
\begin{itemize}
    \item 十进制D：基数10，逢十进一，数码0~9，位权$10^i$
    \item 二进制B：基数2，逢二进一，数码0/1，数字电路基础，位权$2^i$
    \item 八进制O：基数8，逢八进一，数码0~7，位权$8^i$
    \item 十六进制H：基数16，逢十六进一，数码0~9、A(10)~F(15)，位权$16^i$
\end{itemize}

\subsection{0~15多进制对照表}
\begin{table}[!htbp]
    \centering
    \setlength{\tabcolsep}{10pt}
    \renewcommand{\arraystretch}{1.3}
    \caption{0~15十进制、二进制、八进制、十六进制对应关系}
    \label{tab:num_0_15}
    \begin{tabular}{|c|c|c|c|c|c|c|c|}
        \hline
        D & B & O & H & D & B & O & H \\
        \hline
        0  & 0000 & 00 & 0 & 8  & 1000 & 10 & 8 \\
        1  & 0001 & 01 & 1 & 9  & 1001 & 11 & 9 \\
        2  & 0010 & 02 & 2 & 10 & 1010 & 12 & A \\
        3  & 0011 & 03 & 3 & 11 & 1011 & 13 & B \\
        4  & 0100 & 04 & 4 & 12 & 1100 & 14 & C \\
        5  & 0101 & 05 & 5 & 13 & 1101 & 15 & D \\
        6  & 0110 & 06 & 6 & 14 & 1110 & 16 & E \\
        7  & 0111 & 07 & 7 & 15 & 1111 & 17 & F \\
        \hline
    \end{tabular}
\end{table}

\section{不同数制间的转换}
\subsection{任意进制转十进制}
方法：按位权展开求和，核心公式：
\[
(X)_N=\sum_{i=-m}^{n-1}k_i N^i
\]

\subsection{十进制转N进制}
\begin{itemize}
    \item 整数部分：\textbf{除N取余，逆序排列}，直至商为0
    \item 小数部分：\textbf{乘N取整，顺序排列}，至小数为0或满足精度
    \item 特例：转二进制替换基数为2；转八/十六进制同理替换基数
\end{itemize}

\subsection{二进制与八/十六进制互转}
\begin{table}[!htbp]
    \centering
    \setlength{\tabcolsep}{12pt}
    \renewcommand{\arraystretch}{1.3}
    \caption{分组转换规则}
    \label{tab:bin_group_rule}
    \begin{tabular}{|c|c|}
        \hline
        转换方向 & 操作步骤 \\
        \hline
        B $\to$ O & 小数点向两侧每3位分组，不足补0，每组对应1位八进制 \\
        B $\to$ H & 小数点向两侧每4位分组，不足补0，每组对应1位十六进制 \\
        O $\to$ B & 1位八进制展开为3位二进制，删除首尾多余前导0 \\
        H $\to$ B & 1位十六进制展开为4位二进制，删除首尾多余前导0 \\
        \hline
    \end{tabular}
\end{table}

\section{二进制算数运算}
\subsection{基础四则运算规则}
二进制加减乘除遵循：逢二进一、借一当二；乘法/除法可拆解为移位+加法，减法统一转化为\textbf{加补码}。

\subsection{原码、反码、补码}
带符号二进制数由符号位+数值位构成：最高位为符号位，0表示正数，1表示负数，剩余低位为数值位。
\subsubsection{原码}
定义：符号位加上数值绝对值对应的二进制码。
\begin{itemize}
    \item 正数原码：符号位为0，数值位等于真值；
    \item 负数原码：符号位为1，数值位等于真值绝对值；
    \item 特性：数字0存在两种编码，$[+0]_{\text{原}}$、$[-0]_{\text{原}}$，不便于减法运算。
    \item 用途：二进制负数直观表示，仅用于人工阅读，硬件不做运算。
\end{itemize}

\subsubsection{反码}
定义：基于原码按位取反得到。
\begin{itemize}
    \item 正数反码：与自身原码完全相等；
    \item 负数反码：符号位保持1不变，数值位全部按位取反（0变1，1变0）；
    \item 特性：数字0仍存在两种编码，不能统一0的表示。
    \item 用途：原码转换补码的中间过渡编码，无独立运算价值。
\end{itemize}

\subsubsection{补码}
定义：计算机系统实际存储、运算使用的编码，消除0的二义性，实现减法转加法。
\begin{itemize}
    \item 正数补码：补码 = 原码 = 反码；
    \item 负数补码：对负数的反码加1；
    \item $n$位带符号补码取值范围：$\boldsymbol{-2^{n-1} \le X \le 2^{n-1}-1}$；
    \item 特性：数字0仅有唯一补码表示；符号位可直接参与加法运算。
    \item 运算核心公式：$[A-B]_{\text{补}} = [A]_{\text{补}} + [-B]_{\text{补}}$；
    \item 溢出规则：加法产生最高进位直接舍弃；运算结果超出取值范围则发生溢出错误。
    \item 用途：数字电路、CPU唯一运算编码，仅用加法器即可完成加减运算，简化硬件电路；同等位数下可表示更多负数。
\end{itemize}

\subsubsection{4位带符号数原/反/补码对照表}
\begin{table}[!htbp]
    \centering
    \setlength{\tabcolsep}{6pt}
    \renewcommand{\arraystretch}{1.2}
    \caption{4位有符号数编码对照}
    \label{tab:4bit_code}
    \begin{tabular}{|c|c|c|c|}
        \hline
        十进制数 & 原码 & 反码 & 补码 \\
        \hline
        $+7$ & 0111 & 0111 & 0111 \\
        $+6$ & 0110 & 0110 & 0110 \\
        $+5$ & 0101 & 0101 & 0101 \\
        $+4$ & 0100 & 0100 & 0100 \\
        $+3$ & 0011 & 0011 & 0011 \\
        $+2$ & 0010 & 0010 & 0010 \\
        $+1$ & 0001 & 0001 & 0001 \\
        $0$   & 0000 & 0000 & 0000 \\
        $-1$  & 1001 & 1110 & 1111 \\
        $-2$  & 1010 & 1101 & 1110 \\
        $-3$  & 1011 & 1100 & 1101 \\
        $-4$  & 1100 & 1011 & 1100 \\
        $-5$  & 1101 & 1010 & 1011 \\
        $-6$  & 1110 & 1001 & 1010 \\
        $-7$  & 1111 & 1000 & 1001 \\
        \hline
    \end{tabular}
\end{table}

\section{几种常用的编码}
\subsection{十进制代码}
\begin{table}[!htbp]
    \centering
    \setlength{\tabcolsep}{5pt}
    \renewcommand{\arraystretch}{1.2}
    \caption{几种常用十进制代码对照表}
    \label{tab:bcd_all}
    \begin{tabular}{|c|c|c|c|c|c|}
        \hline
        十进制数 & 8421码 & 余3码 & 2421码 & 5211码 & 余3循环码 \\
        \hline
        0 & 0000 & 0011 & 0000 & 0000 & 0010 \\
        1 & 0001 & 0100 & 0001 & 0001 & 0110 \\
        2 & 0010 & 0101 & 0010 & 0100 & 0111 \\
        3 & 0011 & 0110 & 0011 & 0101 & 0101 \\
        4 & 0100 & 0111 & 0100 & 0111 & 0100 \\
        5 & 0101 & 1000 & 1011 & 1000 & 1100 \\
        6 & 0110 & 1001 & 1100 & 1001 & 1101 \\
        7 & 0111 & 1010 & 1101 & 1100 & 1111 \\
        8 & 1000 & 1011 & 1110 & 1101 & 1110 \\
        9 & 1001 & 1100 & 1111 & 1111 & 1010 \\
        \hline
    \end{tabular}
\end{table}
\begin{enumerate}
    \item 8421码：又称BCD码，是最常用的十进制编码。每位从左至右权值依次为8、4、2、1，展开公式 $D=\sum k_i 2^i$。
    \textbf{转换方法：将4位二进制码每一位乘以对应权值，全部相加得到对应十进制数。}
    例：$(1001)_{\text{8421BCD}}=1\times 2^3+0\times 2^2+0\times 2^1+1\times 2^0=8+0+0+1=9$；
    例：$(0110)_{\text{8421BCD}}=0\times 8+1\times 4+1\times 2+0\times 1=6$。
    仅0000~1001有效，1010~1111无对应十进制，属于非法编码。

    \item 余3码：无权码，编码二进制展开数值比对应十进制大3；例0101代表十进制2，二进制展开值为5。做加法和为10时，自动产生进位信号。

    \item 2421码：有权码，权值从左至右2、4、2、1；例$(1100)_2=1\times 2+1\times 4+0\times 2+0\times 1=6$。0与9、1与8、2与7互为反码，两编码相加为9时，结果为1111。

    \item 5211码：有权码，权值从左至右5、2、1、1；例$(0111)_2=0\times 5+1\times 2+1\times 1+1\times 1=4$，多用于分频器。

    \item 余3循环码：无权码，相邻编码仅一位状态不同，译码不会出错。
\end{enumerate}

\subsection{格雷码}
\subsubsection{基本特点与应用}
\begin{itemize}
    \item 每位状态按固定顺序循环；编码顺序递增时，相邻代码仅有1位状态改变。
    \item 作用：减少信号过渡噪声，抗干扰能力强，主要用于计数器、ADC。
\end{itemize}

\begin{table}[!htbp]
    \centering
    \setlength{\tabcolsep}{6pt}
    \renewcommand{\arraystretch}{1.2}
    \caption{4位二进制码与格雷码对照表}
    \label{tab:gray_4bit}
    \begin{tabular}{|c|c|c||c|c|c|}
        \hline
        编码顺序 & 二进制码 & 格雷码 & 编码顺序 & 二进制码 & 格雷码 \\
        \hline
        0 & 0000 & 0000 & 8  & 1000 & 1100 \\
        1 & 0001 & 0001 & 9  & 1001 & 1101 \\
        2 & 0010 & 0011 & 10 & 1010 & 1111 \\
        3 & 0011 & 0010 & 11 & 1011 & 1110 \\
        4 & 0100 & 0110 & 12 & 1100 & 1010 \\
        5 & 0101 & 0111 & 13 & 1101 & 1011 \\
        6 & 0110 & 0101 & 14 & 1110 & 1001 \\
        7 & 0111 & 0100 & 15 & 1111 & 1000 \\
        \hline
    \end{tabular}
\end{table}

\subsubsection{自然码与循环码定义}
\begin{itemize}
    \item 自然码：有权码，每位固定权值，结构同二进制；$n$位最大计数值 $2^n-1$。
    \item 循环码（格雷码）：无权码，各位独立循环；任意相邻码组仅1位不同。
\end{itemize}

\subsection{美国信息交换标准代码（ASCII）}
ASCII为7位二进制代码，总计128个字符，应用于计算机与通讯领域。
\begin{itemize}
    \item 字符'0'：$48_{\text{D}}=0110000_{\text{B}}$
    \item 字符'A'：$65_{\text{D}}=1000001_{\text{B}}$
    \item 字符'a'：$97_{\text{D}}=1100001_{\text{B}}$
\end{itemize}